\thispagestyle{empty}

\begin{center}
{\Large {\bf \uppercase{ABSTRACT}}}
\end{center}

\vspace{\baselineskip}
\hrule
\vspace{2\baselineskip}

\justifying
\setlength{\parskip}{0.6em}

Question--answer based story generation is an emerging direction in
Generative Artificial Intelligence that transforms passive narrative
consumption into active human--AI co-creation. With remarkable
advancements in Large Language Models (LLMs) such as GPT-3, GPT-4,
LLaMA, Mistral, BART and T5, machines have achieved unprecedented
capabilities in understanding natural-language queries, capturing
nuanced user intent and producing personalised, emotionally engaging
stories.

This dissertation presents the design, implementation and evaluation
of a comprehensive Q\&A based interactive storytelling system. The
system uses a structured 17-question elicitation protocol across four
narrative dimensions (setting, character, theme, plot) to gather user
preferences, and routes them through a seven-module backend pipeline
comprising User Input Acquisition, Context Extraction, Knowledge and
Memory, Prompt Engineering, LLM-based Story Generation, Story Flow
and Consistency Management, and Output Formatting. The framework is
provider-agnostic and supports local (Ollama) as well as cloud
inference (Groq, Hugging Face, OpenAI), enabling deployment across a
wide range of cost and latency budgets.

Beyond text, the system delivers a multimodal reader experience:
per-segment scene illustrations are synthesised on demand via the
Pollinations.ai diffusion service using a sentence-scored prompt
extractor and a serial-throttled request queue; spoken narration is
produced through the W3C Web Speech Synthesis interface; and
genre-conditioned ambient music is streamed from a royalty-free CDN
with time-aligned cross-fades. Reader analytics are surfaced through
a D3-based Story Map and a Character Relationship Graph derived from
filtered co-appearance statistics. An end-user authoring layer
admits free-form custom answers per question and entirely user-added
Q\&A pairs, which are routed through a dedicated
\texttt{custom\_elements} context channel into the LLM prompt.

Empirical evaluation across 50 generated stories and 15 user-study
participants reports mean Likert scores of 4.05/5 (coherence),
4.12/5 (creativity), 4.35/5 (engagement) and 3.75/5 (consistency); a
1.8\,s Groq--LLaMA-3.1-8B cloud latency per initial segment; an
overall satisfaction of 4.3/5; and a character-detection $F_1$ of
0.80 after stop-word filtering and a min-2 occurrence threshold.

\vspace{0.5\baselineskip}
{\large \bf{Keywords}: }
Generative AI, Large Language Models, Interactive Storytelling,
Question--Answer Systems, Narrative Generation, Prompt Engineering,
Multimodal Co-creation, Text-to-Image Synthesis, Text-to-Speech,
Human--AI Collaboration, FastAPI, D3.js visualisation.
